% Ad Atticum 9.1 --- headnote
% ad-atticum-09-01, 6 March 49 BC, Formian villa

\textit{Cicero to Atticus, written from the Formian villa
on the day before the Nones of March 49 BC (the manuscript
dateline: \textit{Scr.\ in Formiano prid.\ Non.\ Mart.\ a.\
705 (49)}). The opening letter of book 9. It is the
fourteenth day since Pompey moved from Canusium toward
Brundisium; Cicero, still on the coast, is calculating
hours, watching the road for couriers, and finding the
silence from Brundisium \textit{mirum}. The City has
emptied: the optimates are streaming south, even Manius
Lepidus, the friend with whom he used to ``wear away the
day,'' is going tomorrow.}

\textit{Section 1 sets the count of days and the
unbearable expectation. Section 2 is the search for
Lentulus and Domitius --- where they are, by what road
they will go to Pompey, when --- and the report that
Rome is now full of the leading men converging south.
Section 3 lays out Cicero's own intended itinerary
(Formiae, Arpinum, then by the least-travelled route to
the Adriatic, lictors dismissed) and turns suddenly
bitter: the good men in their well-timed dinner-parties
disapprove of his lingering; very well, let us yield,
make war on Italy by land and sea, light up again the
hatreds of the wicked, follow the counsels of Lucceius
and Theophanes. Section 4 is the roll-call of who is
going and why --- Scipio by allotment or as Pompey's
father-in-law's man or as Caesar's fugitive; the
Marcelli held only by fear of Caesar's sword; Appius
under the same dread, with fresh enmities besides; the
legates and Faustus the proquaestor going on duty;
Cicero alone free to choose either way. Quintus comes
too, whom it was not right to entangle in this fortune
of his brother's. The closing sentence is bare: this
concession is made to Pompey alone, who does not even
ask for it, claiming to plead not his own cause but the
commonwealth's. ``I should very much like to know what
you are thinking about crossing over into Epirus.''}
